\documentclass[11pt,a4paper]{article}

\usepackage[margin=2.5cm]{geometry}
\usepackage{amsmath, amssymb}
\usepackage{booktabs}
\usepackage{xcolor}
\usepackage{hyperref}
\usepackage{listings}
\usepackage{graphicx}
\usepackage{caption}
\usepackage{subcaption}
\usepackage{parskip}

% Colors
\definecolor{darkgreen}{HTML}{1A3A2A}
\definecolor{gold}{HTML}{D4A843}
\definecolor{codeblue}{HTML}{2E4A7A}
\definecolor{codegray}{HTML}{555555}

% Code listing style
\lstset{
  basicstyle=\ttfamily\small,
  keywordstyle=\color{codeblue}\bfseries,
  commentstyle=\color{codegray}\itshape,
  breaklines=true,
  frame=single,
  backgroundcolor=\color{gray!8},
  numbers=left,
  numberstyle=\tiny\color{codegray},
  language=Python,
}

\hypersetup{colorlinks=true, linkcolor=codeblue, urlcolor=codeblue}

% ─────────────────────────────────────────────────────────────────────────────

\title{
  \textbf{LLM-Guided Evolutionary Search} \\[0.4em]
  \large MSc Thesis Progress Report \\[0.2em]
  \normalsize Technion --- Israel Institute of Technology
}
\author{Shira Ahtan \\ \texttt{sgahtan@nvidia.com}}
\date{September 2026}

\begin{document}
\maketitle
\tableofcontents
\newpage

% ─────────────────────────────────────────────────────────────────────────────
\section{Motivation and Thesis Statement}

\subsection{Background: AlphaEvolve}

AlphaEvolve~\cite{alphaevolve2025} (Google DeepMind, 2025) is an evolutionary
program-search framework that uses large language models (LLMs) as mutation
operators. Starting from a seed program, AlphaEvolve iteratively proposes
improved variants using Gemini, evaluated on a fixed benchmark. Its most
notable result was discovering a matrix multiplication algorithm that beat
Strassen's 56-year-old record: reducing complexity from $O(n^{2.807})$ to
$O(n^{2.778})$.

AlphaEvolve's mutation step can be written as:
\[
  x_{t+1} = \mathrm{LLM}(x_t)
\]

The LLM sees only the current code $x_t$. It is \textbf{blind to why the
program failed}: it does not receive the score, the error trace, or any
signal about what went wrong.

\subsection{The Thesis Contribution}

This thesis proposes one targeted upgrade: provide the LLM with \emph{failure
evidence} $E(x_t)$ before each mutation:

\[
  \boxed{x_{t+1} = \mathrm{LLM}\!\left(x_t,\; E(x_t),\; \mathrm{score}(x_t)\right)}
\]

where $E(x_t)$ is a structured failure report containing:
\begin{itemize}
  \item The current score (0--100)
  \item Correctness: how many test cases passed
  \item Speed: measured speedup vs baseline
  \item The last exception message (if any)
  \item A truncated stack trace (if any)
\end{itemize}

The hypothesis is that this \textbf{diagnosis-driven mutation} --- which we
call the \emph{smart step} --- should reach a target fitness with fewer
evaluator calls than the blind mutation (\emph{dumb step}), improving
\emph{sample efficiency}.

We call the blind baseline the \textit{dumb step}:
\[
  x_{t+1} = \mathrm{Mutate\_random}(x_t)
\]

% ─────────────────────────────────────────────────────────────────────────────
\section{Research Questions}

\begin{enumerate}
  \item[\textbf{Q1}] \textbf{Sample efficiency}: Does failure-informed mutation
    reach target fitness with fewer evaluator calls?
  \item[\textbf{Q2}] \textbf{Explore vs exploit}: When in the search trajectory
    does the smart step help most --- early exploration or late refinement?
  \item[\textbf{Q3}] \textbf{Smart/dumb balance}: What is the optimal ratio of
    smart to dumb steps, especially when diversity collapses?
  \item[\textbf{Q4}] \textbf{LLM convergence}: Does using the same LLM for all
    mutations cause the population to converge to homogeneous code patterns?
  \item[\textbf{Q5}] \textbf{Failure signal}: Which component of $E(x_t)$ is
    most informative --- error type, score delta, or stack trace?
\end{enumerate}

% ─────────────────────────────────────────────────────────────────────────────
\section{System Architecture}

The implementation is modular, with each concern in its own file:

\begin{center}
\begin{tabular}{lp{9cm}}
\toprule
\textbf{File} & \textbf{Responsibility} \\
\midrule
\texttt{config.py}          & All hyperparameters in one place \\
\texttt{benchmark.py}       & Fixed sorting task; evaluator returning score 0--100 \\
\texttt{benchmark\_matrix.py} & Second task: matrix multiplication \\
\texttt{population.py}      & Population of programs; weighted parent selection \\
\texttt{mutator.py}         & Smart step (LLM) and dumb step (random + structural) \\
\texttt{evaluators.py}      & Multiple scoring variants (combined, correctness, speed) \\
\texttt{scheduler.py}       & Decides smart vs dumb per generation \\
\texttt{loop.py}            & Main evolutionary loop; saves JSON logs \\
\texttt{main.py}            & CLI entry point \\
\texttt{compare.py}         & Ablation runner: smart vs dumb vs mixed \\
\texttt{param\_sweep.py}    & Hyperparameter grid search (no LLM required) \\
\texttt{plot.py}            & Three-panel results visualization \\
\bottomrule
\end{tabular}
\end{center}

\subsection{The Evolutionary Loop}

Each generation:
\begin{enumerate}
  \item Select a parent from the top-$K$ programs (weighted by score)
  \item Decide: smart step or dumb step (via Scheduler)
  \item Generate $k$ candidate mutations in parallel (multiple proposers)
  \item Evaluate all candidates on the fixed benchmark
  \item Accept the best candidate if it beats the worst in the population
  \item Log and optionally early-stop if target score is reached
\end{enumerate}

\subsection{Dumb Step Mutations}

The dumb step applies one of eight random mutations without any knowledge of
why the program failed:

\begin{enumerate}
  \item Swap two adjacent lines
  \item Replace a numeric literal with a nearby value ($\pm 1, \pm 2$)
  \item Swap a comparison or arithmetic operator
  \item Duplicate a line
  \item Delete a non-essential line
  \item \textbf{Insert a sorting primitive} (insertion sort, merge sort, quicksort, heapsort)
  \item \textbf{Add an early-exit guard} (\texttt{if len(arr) <= 1: return})
  \item \textbf{Rename a loop variable} (tests population diversity)
\end{enumerate}

Items 6--8 are \emph{structural mutations} added to make the dumb-step
baseline stronger --- a better baseline makes the comparison with the smart
step more meaningful.

\subsection{Evaluator Variants}

Four scoring functions are implemented in \texttt{evaluators.py}:

\begin{center}
\begin{tabular}{lp{8cm}}
\toprule
\textbf{Variant} & \textbf{Formula} \\
\midrule
\texttt{combined}      & $70\% \times \text{correctness} + 30\% \times \text{speed}$ (default) \\
\texttt{correctness}   & $100\% \times \text{correctness}$, speed ignored \\
\texttt{speed}         & $100\% \times \text{speed}$, hard gate on correctness \\
\texttt{strict\_speed} & $\min(100,\; 20 \times \text{speedup})$, finer resolution \\
\bottomrule
\end{tabular}
\end{center}

% ─────────────────────────────────────────────────────────────────────────────
\section{Benchmark Tasks}

\subsection{Sorting (Primary)}

\textbf{Task}: Optimize a \texttt{sort\_array(arr: list) -> list} function.

\textbf{Seed program}: Selection sort ($O(n^2)$).

\textbf{Scoring}:
\[
  \text{score} = 70 \cdot \frac{\text{correct}}{\text{total}} + \min\!\left(30,\; 30 \cdot \frac{t_{\text{baseline}}}{t_{\text{candidate}}}\right)
\]

\textbf{Test cases}: 20 fixed arrays (seed 42), sizes 10--200, values
$[-1000, 1000]$.

\textbf{Why sorting?} It is deterministic, fully reproducible, gradable on
two axes (correctness + speed), and directly analogous to AlphaEvolve's
matrix multiplication task.

\subsection{Matrix Multiplication (Second Benchmark)}

\textbf{Task}: Optimize a \texttt{matmul(A, B) -> list} function.

\textbf{Seed program}: Naive $O(n^3)$ triple-loop multiplication.

\textbf{Scoring}: $60\%$ correctness $+$ $40\%$ speed.

\textbf{Test cases}: 15 cases ($5 \times$ each of $2\times2$, $4\times4$, $8\times8$
matrices).

\textbf{Why matrix multiplication?} This is exactly the task where AlphaEvolve
made its headline result. Running our framework on this task directly compares
to the AlphaEvolve baseline.

% ─────────────────────────────────────────────────────────────────────────────
\section{Preliminary Results}

\subsection{Dumb-Only Baseline (6 Seeds)}

We ran the dumb-only configuration (equivalent to AlphaEvolve's blind mutation)
across 6 random seeds to establish a robust baseline.

\begin{center}
\begin{tabular}{cccc}
\toprule
\textbf{Seed} & \textbf{Start Score} & \textbf{Final Score} & \textbf{Target (95) Reached?} \\
\midrule
42  & 70.4 & 71.6 & \texttimes \\
1   & 70.4 & 72.2 & \texttimes \\
7   & 70.4 & 70.5 & \texttimes \\
13  & 70.4 & 70.5 & \texttimes \\
21  & 70.4 & 70.5 & \texttimes \\
99  & 70.4 & 70.5 & \texttimes \\
\midrule
\textbf{Mean} & \textbf{70.4} & \textbf{71.1} & \textbf{0/6} \\
\bottomrule
\end{tabular}
\end{center}

\textbf{Key observations}:
\begin{itemize}
  \item Blind random mutation is unable to escape the selection sort basin.
    All 6 runs plateau at 70--73 within 15 generations.
  \item Population diversity collapses to 0.0 by generation $\approx$15 in
    every run. Once everyone in the population is the same program, dumb steps
    produce only broken variants that fail to score higher.
  \item The ``best'' program found by dumb-only is selection sort with minor
    cosmetic edits (e.g.\ a duplicated \texttt{return} line).
  \item All 251 evaluations are consumed (50 generations $\times$ 5 proposers
    $+$ 1 seed), yet the target score of 95 is never reached.
\end{itemize}

\textbf{Interpretation}: This is exactly the motivation for the smart step.
The LLM, given the failure evidence (``speed: 0.02$\times$ baseline, correct
but $O(n^2)$''), should be able to propose a fundamentally different algorithm
--- something a random line swap cannot do.

\subsection{Evaluator Sanity Check}

Running all four evaluator variants on the selection sort seed program:

\begin{center}
\begin{tabular}{lccr}
\toprule
\textbf{Evaluator} & \textbf{Score} & \textbf{Correct} & \textbf{Speedup} \\
\midrule
\texttt{combined}     & 70.6 & 20/20 & 0.02$\times$ \\
\texttt{correctness}  & 100.0 & 20/20 & --- \\
\texttt{speed}        & 2.4  & 20/20 & 0.02$\times$ \\
\texttt{strict\_speed} & 0.5  & 20/20 & 0.02$\times$ \\
\bottomrule
\end{tabular}
\end{center}

The \texttt{correctness} evaluator gives 100/100 immediately (selection sort
is correct), making this a trivial task. The \texttt{speed} and
\texttt{strict\_speed} evaluators give nearly 0 (selection sort is 50$\times$
slower than Python's built-in Timsort), making the search objective pure
speed optimization. This validates all four variants behave as intended.

\subsection{Matrix Benchmark Sanity Check}

The naive $O(n^3)$ seed program on the matrix task:
\begin{itemize}
  \item Score: 97.8/100
  \item Correct: 15/15 test cases
  \item Speedup: 0.95$\times$ vs baseline (essentially identical --- baseline \emph{is} naive)
\end{itemize}

The high score means the matrix task is less challenging at the baseline level.
The interesting question is whether the search can discover Strassen-like
divide-and-conquer approaches that go \emph{faster} than naive.

% ─────────────────────────────────────────────────────────────────────────────
\section{Next Steps}

\subsection{Immediate (requires API key)}
\begin{enumerate}
  \item Set \texttt{ANTHROPIC\_API\_KEY} and run \texttt{python compare.py}:
    compare smart-only vs dumb-only vs mixed across 3+ seeds.
  \item Analyse the smart step trajectories: does the LLM propose merge sort /
    quicksort on the first few calls?
  \item Run the evaluator ablation: does \texttt{correctness} or \texttt{speed}
    as the sole signal change what the LLM proposes?
\end{enumerate}

\subsection{Short Term}
\begin{enumerate}
  \item Complete the hyperparameter sweep (\texttt{param\_sweep.py}) and
    identify the best population/proposer configuration.
  \item Port the framework to the matrix multiplication benchmark and run
    the same smart vs dumb comparison.
  \item Add AutoResearch~\cite{autoresearch} as a third comparison baseline.
\end{enumerate}

\subsection{Long Term}
\begin{enumerate}
  \item Scale to 3+ seeds per configuration and report mean $\pm$ std.
  \item Study diversity collapse: measure code-level similarity (token overlap,
    AST edit distance) not just score std-dev.
  \item Investigate island model: run $k$ independent populations and
    periodically migrate the best programs between them.
  \item Ablate $E(x_t)$ components (Q5): does removing the stack trace or
    the score from the context change the LLM's proposals?
\end{enumerate}

% ─────────────────────────────────────────────────────────────────────────────
\section{Repository}

Code is available at:
\url{https://github.com/shiragahtan/thesis_deep}

\begin{thebibliography}{9}
\bibitem{alphaevolve2025}
  Google DeepMind,
  \textit{AlphaEvolve: A Gemini-Powered Coding Agent for Designing Advanced Algorithms},
  2025.

\bibitem{funsearch2023}
  Romera-Paredes et al.,
  \textit{Mathematical Discoveries from Program Search with Large Language Models (FunSearch)},
  Nature, 2023.

\bibitem{autoresearch}
  AutoResearch: Automated Scientific Research using LLM-Guided Search,
  2024.
\end{thebibliography}

\end{document}
